\chapter{Introduzione}
Nel panorama attuale della sicurezza informatica, le botnet rappresentano una delle minacce più rilevanti e persistenti per le infrastrutture digitali moderne, risultando particolarmente difficili da individuare e smantellare in maniera definitiva. La loro natura distribuita e dinamica consente infatti agli attaccanti di mascherare l'origine delle operazioni e di adattare rapidamente il comportamento della rete malevola in risposta ad eventuali interventi di mitigazione. Queste reti malevole sono costituite da una moltitudine di dispositivi compromessi, denominati bot o zombie, che operano sotto il controllo di uno o più attaccanti attraverso infrastrutture di comando e controllo. Gli attaccanti sfruttano queste reti per svolgere una vasta gamma di attività illecite, beneficiando dell'elevato numero di dispositivi coinvolti. Tra le principali attività rientrano attacchi Denial-of-Service (DoS) e Distributed Denial-of-Service (DDoS) su larga scala, in grado di rendere servizi e applicazioni temporaneamente inaccessibili. A questi si affiancano l'invio massivo di malware e spam, finalizzato alla compromissione di ulteriori dispositivi e all'espansione della botnet stessa, nonché campagne di phishing e frodi informatiche mirate alla raccolta di informazioni sensibili, come credenziali di accesso e dati personali. In molti casi, le botnet vengono impiegate come infrastrutture a noleggio, rendendole strumenti estremamente insidiosi e utilizzabili per scopi differenti a seconda delle esigenze degli attaccanti. 

L'evoluzione delle infrastrutture di rete e la crescita esponenziale del numero di dispositivi connessi hanno contribuito in modo significativo ad accrescere la portata di questo fenomeno. In particolare, la diffusione dei dispositivi Internet of Things (IoT) ha introdotto nuove superfici di attacco, spesso caratterizzate da limitate capacità di difesa e da configurazioni non adeguate. Dispositivi come telecamere IP, router domestici, videocitofoni, smart TV e altri elettrodomestici connessi alla vengono distribuiti con credenziali di default, firmware obsoleti o meccanismi di aggiornamento assenti, rendendoli bersagli ideali per infezioni automatiche su larga scala. 

La nascita e l'evoluzione delle botnet IoT evidenziano come la sicurezza informatica non riguardi più esclusivamente i computer tradizionali, ma coinvolga qualsiasi dispositivo dotato di connettività di rete. La compromissione di apparati completamente innocui può infatti avere conseguenze rilevanti sulla stabilità e sulla disponibilità di servizi critici, con impatti che si estendono ben oltre il singolo dispositivo infetto e che possono interessare intere infrastrutture di rete. In questo contesto si colloca il malware Mirai, responsabile della creazione di una delle botnet IoT più conosciute e studiate, che ha attirato l'attenzione della comunità scientifica e industriale per la semplicità delle sue tecniche di infezione e per l'elevata insidiosità dei suoi attacchi. Mirai ha dimostrato come lo sfruttamento di vulnerabilità banali, quali l'utilizzo di credenziali deboli o predefinite, possa tradursi in una capacità offensiva estremamente elevata, capace di generare volumi di traffico tali da compromettere la disponibilità di importanti servizi di rete. 

Lo studio di una botnet come Mirai risulta particolarmente rilevante, in quanto rappresenta un caso emblematico in cui la semplicità di infezione dei dispositivi vulnerabili si combina con un'architettura di controllo efficiente e con un alto potenziale di attacco. Analizzare Mirai consente di comprendere non solo le tecniche di propagazione del malware, ma anche le strategie di gestione dei bot, e i meccanismi con cui questi attacchi possono impattare in maniera significativa su server e infrastrutture di rete.


Alla luce delle precedenti considerazioni, la presente tesi si propone di esaminare in modo approfondito la botnet Mirai attraverso la realizzazione di un ambiente di simulazione controllato. La scelta di focalizzarsi su Mirai non è motivata solo dalla sua rilevanza storica, ma anche dalla sua persistente attualità nel panorama attuale. Infatti nonostante la pubblicazione del codice sorgente e la conseguente diffusione di numerose varianti, Mirai continua a rappresentare un modello di riferimento per lo studio delle botnet IoT. 

L'analisi di Mirai consente di evidenziare come vulnerabilità semplici, spesso trascurate in fase di progettazione, possano essere sfruttate in modo sistematico per costruire infrastrutture di attacco estremamente pericolose. Comprendere tali dinamiche risulta fondamentale sia dal punto di vista offensivo sia da quello difensivo, poiché permette di individuare le debolezze più critiche nei dispositivi IoT e di valutare l'efficacia delle contromisure adottabili. In questo senso, lo studio di Mirai offre una base concreta per la comprensione di fenomeni più complessi e per l'analisi di botnet di nuova generazione. 

L'utilizzo di un ambiente di test controllato consente di studiare in maniera sistematica tutte le fasi operative del malware, dalla scansione della rete alla ricerca dei dispositivi vulnerabili, al processo di infezione e di integrazione degli stessi all'interno della botnet, fino all'esecuzione degli attacchi DoS/DDoS, analizzando il traffico generato e l'impatto sulla disponibilità del server bersaglio.

Infine l'attenzione sarà dedicata all'analisi degli attacchi DoS/DDoS per comprendere al meglio uno dei principali vettori di minaccia per la disponibilità dei servizi di rete. Questi attacchi, sempre più frequenti e sofisticati, rappresentano una minaccia critica per le organizzazioni, rendendo fondamentale lo studio dei meccanismi attraverso cui essi vengono realizzati e degli effetti che producono sui sistemi bersaglio. 

L'obiettivo è quello di fornire una visione chiara e dettagliata del funzionamento di Mirai in un contesto controllato, mettendo in evidenza sia le criticità di sicurezza che rendono possibili tali attacchi, sia gli effetti concreti che essi possono avere sui servizi di rete e sulle risorse dei sistemi bersaglio. In particolare il contenuto della tesi è organizzato come segue:
\begin{itemize}
    \item \textbf{Capitolo 1}: fornisce una presentazione generale degli argomenti trattati nella tesi;
    \item \textbf{Capitolo 2}: illustra il background teorico definendo i concetti fondamentali necessari a una corretta comprensione del lavoro, tra cui quello di botnet e di attacco DoS/DDoS. Viene inoltre analizzato il malware Mirai con particolare attenzione alle sue modalità di propagazione, alla struttura della botnet conseguente all'infezione dei dispositivi vulnerabili e alle capacità offensive della stessa;
    \item \textbf{Capitolo 3}: descrive in dettaglio tutti gli strumenti e le tecnologie impiegati per la creazione dell'ambiente di simulazione. Vengono anche illustrati gli strumenti per il monitoraggio della rete e delle risorse del server bersaglio;
    \item \textbf{Capitolo 4}: espone i risultati ottenuti durante la creazione della botnet e l’esecuzione degli attacchi verso il server bersaglio, accompagnati da un’analisi dettagliata del traffico generato nella rete e degli effetti osservati sulla disponibilità del servizio;
    \item \textbf{Capitolo 5}: presenta le conclusioni tratte dai risultati ottenuti, vengono inoltre illustrate delle possibili estensioni del lavoro realizzato.
\end{itemize}



